\section*{A first pass: what the proof has to do}\label{sec:first-pass}
\addcontentsline{toc}{section}{A first pass: what the proof has to do}
An upper bound limits the size of \emph{every} Condorcet domain.
An explicit construction supplies one domain and hence a lower bound.
We will use small constructions to make the counting visible, and
then explain why the same counting limits all domains.

\subsection*{1. Erase a label and count what survives}
Suppose a domain \(D\) has \(M\) orders on \(n\) alternatives.
Erasing a label can merge different orders. The proof asks how much
merging is possible, averaged over the \(n\) choices of erased label.
Here are the four quantities to keep in view:
\begin{center}\small
\begin{tabular}{@{}>{\raggedright\arraybackslash}p{0.08\linewidth}>{\raggedright\arraybackslash}p{0.22\linewidth}>{\raggedright\arraybackslash}p{0.60\linewidth}@{}}\toprule
Symbol & What we count & How to picture it\\\midrule
\(M\) & Orders in \(D\) & The original collection of rankings.\\[0.35em]
\(\Sigma\) & Distinct deletion outputs & An output remembers both the deleted label and the remaining order.\\[0.35em]
\(G\) & Separations in insertion positions & Several separated runs of possible positions for one reinserted label contribute extra units.\\[0.35em]
\(A\) & Enabled adjacent swaps & Each swap edge counts twice, once from either endpoint.\\\bottomrule
\end{tabular}
\end{center}
An insertion fibre is just the list of places where one erased label
could have been. In such a list,
\[
 \text{number of runs}=\text{occupied places}-\text{edges between occupied neighbours}.
\]
Summing this elementary path identity gives
\(\Sigma+G+A=nM\). This is the counting framework for the whole proof.

Every deletion of a Condorcet domain is still Condorcet. Therefore
\(\Sigma\leq nf(n-1)\). To bound \(M\), we need a lower bound
on \(\Sigma\): the orders cannot all disappear into very few outputs.

\subsection*{2. Prove two estimates, for different reasons}
The structural argument gives, for a maximal domain,
\[
 \boxed{G+(n+1)M\leq3\Sigma.}\tag{gap count}
\]
Its main idea is an injective assignment: each freely movable label
forces a distinct obstruction elsewhere. Once that assignment is
proved, the inequality follows by counting moves. No finite certificate
is used in this step.

The second estimate is
\[
 \boxed{\Sigma+G\geq\frac38nM.}\tag{component count}
\]
It comes from splitting a full domain into adjacent-swap components.
Each component is the full solution set of extreme triple laws and
fixed pair comparisons. On these simpler pieces, insertion positions
form intervals and the finite argument proves that deletion retains
at least a \(3/8\) fraction of the \(nM\) order/label incidences.
Splitting counts separated insertion runs separately, which explains
the additional \(G\).

\subsection*{3. The fraction \(32/11\) is now forced}
The two estimates give
\[
 \frac38nM\leq\Sigma+G\leq4\Sigma-(n+1)M.
\]
Move the \((n+1)M\) term to the left and multiply by eight:
\[
 (11n+8)M\leq32\Sigma\leq32nf(n-1).
\]
Choose a domain of maximum size \(M=f(n)\). This is the advertised
recurrence, and its limiting multiplier is \(32/11\).

One can also see the arithmetic without large counts. Set
\(s=\Sigma/(nM)\), \(g=G/(nM)\), with \(n,M>0\). The estimates say
\[
 s+g\geq\frac38,\qquad g+1+\frac1n\leq3s.
\]
Together they force
\(s\geq11/32+1/(4n)\). The term \(1/(4n)\) accounts exactly
for the \(+8\) in the recurrence denominator.
\lean{This overview restates \path{CDCert.TwoInputs.first_input},
\path{second_input} and \path{transfer}. The normalised variables
are an explanatory change of units. The precise unnormalised
statements and their hypotheses appear beside the later proofs.}

\subsection*{4. Where the computer enters, and why copies help}
For the simpler pieces, the difficult input is a particular inequality
on seven alternatives. It bounds a weighted swap count by \(5/8\)
of the family size, plus a correction at each end. The finite
certificate proves this inequality for every permitted constraint
system.

Fixing the order outside a seven-position window lets us apply that
small inequality inside a long order. Consecutive windows share six
positions. Their corrections cancel there, leaving a bounded error
at the two outside ends. Ordered copies then remove this error:
after dividing by the family size, the main counts grow in proportion to the number of copies, while
the total number of outside ends remains two. Thus the exact swap
estimate is \(A\leq5nM/8\). Since interval fibres give
\(\Sigma=nM-A\), the retained deletion fraction is \(3/8\).

This also explains the peak-pit multiplier \(8/3\): applying
\(\Sigma\geq3nM/8\) to a maximum peak-pit domain and using
\(\Sigma\leq nh(n-1)\) gives \(h(n)\leq(8/3)h(n-1)\).
The general-domain proof transfers this \emph{deletion estimate}
through components and the gap count. A geometric estimate for
\(h(n)\) alone would not supply that structural information.

\clearpage
\subsection*{The twelve ideas as reusable mathematical moves}
The following are conceptual analogies, not claims of historical priority.
They also distinguish discovery from the eventual proof.
\begin{center}\small
\renewcommand{\arraystretch}{1.15}
\begin{tabular}{@{}>{\raggedright\arraybackslash}p{0.04\linewidth}>{\raggedright\arraybackslash}p{0.31\linewidth}>{\raggedright\arraybackslash}p{0.54\linewidth}@{}}\toprule
& Move in this proof & General mathematical pattern\\\midrule
1 & Count deletion outputs & Double counting followed by induction.\\
2 & Count runs by places minus edges & The vertices-minus-edges formula for a disjoint union of paths.\\
3 & Assign distinct obstructions & An injection, or a charging argument without reuse.\\
4 & Project using majority & Closure under a median operation and projection in Boolean coordinates.\\
5 & Fix the exterior & Conditioning, then interchanging finite sums.\\
6 & Cancel overlap corrections & Telescoping sums; a discrete difference disappears in the interior.\\
7 & Repeat ordered copies & Amplification: divide a bounded error by an arbitrarily large parameter.\\
8 & Search for useful weights & Linear programming and counterexample-guided discovery.\\
9 & Cover every constraint choice & Exhaustive case analysis with inductively proved soundness.\\
10 & Pay exceptional deficits & Redistribution of local surplus, often called discharging.\\
11 & Check exact finite data & Certificate checking by a small trusted verifier.\\
12 & Relate the two representations & A proved bijection preserving the properties and counts used.\\\bottomrule
\end{tabular}
\end{center}
Idea 8 discovers a candidate; it is not a premise of the theorem.
Ideas 9--11 establish the finite input. Idea 12 ensures that this input
really applies to the preference orders in the original question.
